\section{Equal parameters and the torus--sphere correspondence}
\label{sec:torussphere}

The familiar punctured correspondence extends geometrically to equal
positive boundaries. Let $Q(c^4)$ be a four-holed sphere with all four
boundary lengths $c$, and let $T(2c)$ be a one-holed torus with boundary
length $2c$. The relevant construction is through a common orbifold:
\begin{equation}
 T(2c)\xrightarrow{\;2:1\;}
 O(c;\pi,\pi,\pi)
 \xleftarrow{\;4:1\;}Q(c,c,c,c).
 \label{eq:commonorbifold}
\end{equation}
The middle surface is a sphere with one geodesic boundary and three
angle-$\pi$ orbifold points. The torus quotient uses its elliptic
involution; the sphere quotient uses its Klein-four symmetry. This works
for every $c>0$ and limits to the cusped case~\cite[Section 3]{HuTan}.

A matching simple curve on $T$ covers an arc between two orbifold points
twice; its counterpart on $Q$ covers that arc four times. Thus
\begin{equation}
 \ell_Q=2\ell_T.
 \label{eq:lengthtorussphere}
\end{equation}
The construction gives a correspondence of simple-curve classes, not only
an agreement of dimension or of volume polynomials.

WP forms scale by the degree of a local-isometric covering on a family of
lifted metrics. This follows from naturality of the Goldman symplectic
pairing: its surface integral is multiplied by the covering degree.
On the common marked orbifold family,
\[
 \omega_T=2\omega_O,\qquad \omega_Q=4\omega_O,
 \qquad\text{so}\qquad\omega_Q=2\omega_T.
\]
This comparison takes place before dividing by mapping classes. Branch
points are cone points of the quotient, so the maps are local isometries
in the orbifold sense; their orders $2$ and $4$ multiply the same integral
pairing. The statement does not involve varying the boundary length $c$.
Generically the equal-boundary sphere has its Klein-four isometry group.
There are $4!/4=6$ inequivalent ways to label its boundaries. The torus has
its generic elliptic involution, which acts ineffectively on Teichm\"uller
space. The Klein-four transformations permute the four sphere boundaries
freely and transitively, so none survives after all four labels are fixed.
The generic degree from labeled spheres to the effective torus moduli
space is therefore $6$, and the local area ratio is $2$. Consequently
\begin{equation}
 \boxed{
 V_{0,4}(c,c,c,c)=12V_{1,1}^{\rm eff}(2c)
                =24V_{1,1}^{\rm stack}(2c).}
 \label{eq:voltorussphere}
\end{equation}
Indeed, $24((2c)^2+4\pi^2)/48=2\pi^2+2c^2$, as required by
\eqref{eq:basevolumes}. Exceptional extra automorphisms are accounted for
by orbifold integration. Forgetting the ineffective torus involution is
the source of a common factor-of-two discrepancy.

The cone analogue can be constructed directly. For $0<\phi<\pi$, the
elliptic quotient of a torus with one cone opening $2\phi$ is a sphere
with openings $(\phi,\pi,\pi,\pi)$. Its Klein-four cover branched over
the three angle-$\pi$ points has four openings equal to $\phi$. Uniqueness
of conical uniformization makes the conformal symmetries isometries.
The same degree and labeling argument gives
\begin{equation}
 V_{0,4}(i\phi,i\phi,i\phi,i\phi)
   =2\pi^2-2\phi^2=24V_{1,1}(2i\phi).
 \label{eq:conetorus}
\end{equation}
The common quotient cone surface has area $\pi-\phi$, so $\phi<\pi$ is exactly the
hyperbolic range. There are no collision corrections in this equal-cone
sphere chamber.
The torus has one cone opening $2\phi<2\pi$ and no other noncuspidal
ends, so its polynomial evaluation is covered by the single-cone regime
of Section~\ref{sec:cones}. Thus both sides of~\eqref{eq:conetorus} are
geometric volumes in the stated interval. The endpoint $\phi=\pi$ has
zero area and is a limit, not an additional hyperbolic moduli space.

This correspondence does not assert a direct unbranched degree-two map
from the four-holed sphere to the one-holed torus. The common orbifold and
its two different covering degrees are the precise geometric statement.

